% Chapter 1: The AKS Primality Test
\let\LocalMacro\relax

\chapter{AKS Primality Test}

\section{The Problem}

\subsection{Informal}

\textbf{Background.} A prime number is a whole number greater than 1 that cannot be divided evenly by any smaller whole number except 1 and itself. Testing whether a number is prime is fundamental to cryptography and number theory. For large numbers, the obvious method---trying every possible divisor---is too slow.

\textbf{Given.} You are given a large positive integer $n$.

\textbf{Goal.} Determine whether $n$ is prime or composite, using a procedure that is both fast (runs in polynomial time) and always correct (never gives the wrong answer).

\textbf{Constraints.} The algorithm must be deterministic (no randomness). It must run in polynomial time in the number of digits of $n$. It must give the correct answer for every input. It must not rely on any unproven mathematical assumptions.

\textbf{Example.} Take $n = 97$. Trial division checks divisors 2 through 9, finds none divide 97, and declares it prime. For $n = 104729$ (the 10,000th prime), trial division must check up to about 323---still feasible. But for a 200-digit number, trial division would take longer than the age of the universe. The AKS test, proved correct in 2002, gives the right answer for any number in polynomial time.

\subsection{Formal}
The AKS conjecture (now theorem) asks: is there a deterministic polynomial-time algorithm that correctly decides whether a given integer $n$ is prime or composite, without relying on any unproven assumptions?

\subsection{Historical Context}
The question was motivated by cryptography: RSA encryption relies on the difficulty of factoring, but key generation requires efficient primality testing. Before AKS, cryptographers relied on probabilistic tests.

\subsection{What Was Known}
For centuries, trial division---checking every possible divisor---was the only method. It works but becomes impossibly slow for large numbers. In the 1970s and 1980s, probabilistic tests like Miller--Rabin were developed: they are extremely fast and almost always correct, but there is a tiny chance of error. The question of whether a fast, always-correct test existed was open for decades.

\subsection{Why It Matters}
Testing whether a number is prime is fundamental to cryptography. RSA encryption, which secures online banking and commerce, requires generating huge prime numbers. Before AKS, mathematicians relied on tests that were fast but could occasionally give wrong answers. A deterministic, always-correct, fast test was the holy grail.

\subsection{Simplified Version}
For specific forms of numbers, fast primality tests already existed. Mersenne numbers ($2^p - 1$) can be tested using the Lucas--Lehmer test, and Proth's theorem handles numbers of the form $k \cdot 2^n + 1$. The challenge was a \emph{universal} deterministic test that works for \emph{all} numbers and runs in polynomial time---without relying on any unproven assumptions like the Extended Riemann Hypothesis.

\section{Mathematical Theory}

\subsection{Prerequisites}
This chapter assumes familiarity with:
\begin{itemize}
\item Basic modular arithmetic
\item Polynomial arithmetic over finite fields
\item Complexity theory (P, NP, polynomial time)
\end{itemize}

\subsection{The Key Observation}

Fermat's Little Theorem states that if $p$ is prime and $a \not\equiv 0 \pmod{p}$, then:
\begin{equation}
a^{p-1} \equiv 1 \pmod{p}
\end{equation}

This gives a primality test: check if $a^{n-1} \equiv 1 \pmod{n}$. But Carmichael numbers (like 561) fool this test for all bases coprime to $n$.

The AKS insight was to lift this to \emph{polynomials}:

\begin{theorem}[AKS Criterion]
Let $n > 1$ be an integer and $a$ coprime to $n$. Then $n$ is prime if and only if:
\begin{equation}
(X + a)^n \equiv X^n + a \pmod{n}
\end{equation}
as polynomials in $\mathbb{Z}/n\mathbb{Z}[X]$.
\end{theorem}

This is equivalent to checking that the binomial coefficients $\binom{n}{k}$ are all divisible by $n$ for $1 \leq k \leq n-1$.

\subsection{The Algorithm}

The AKS algorithm (Agrawal, Kayal, Saxena, 2002) proceeds in several steps:

\begin{enumerate}
\item \textbf{Perfect power check}: Is $n = a^b$ for integers $a, b > 1$? If so, composite.
\item \textbf{Find small factor}: Check if $\gcd(a, n) > 1$ for small $a$.
\item \textbf{Size check}: If $n \leq r$, declare prime (where $r$ is a carefully chosen bound).
\item \textbf{Polynomial congruence}: Check $(X + a)^n \equiv X^n + a \pmod{X^r - 1, n}$ for a range of values of $a$.
\end{enumerate}

The key innovation was showing that a \emph{small} value of $r$ suffices, making the algorithm polynomial time.

\section{The Solution}

\subsection{The Big Idea}
There is a beautiful algebraic identity: if $p$ is prime, then $(X + 1)^p$ distributes over addition perfectly---every binomial coefficient $\binom{p}{k}$ for $0 < k < p$ is divisible by $p$. Composite numbers fail this test. That is the Frobenius endomorphism, and it is a *perfect litmus test for primality*. The identity was known for over a century. The breakthrough was proving you can *check it efficiently*: you don't need to expand a polynomial of degree $n$ (which would be impossibly slow). By working modulo a small, cleverly chosen polynomial, you can run the check in polynomial time. It is like having a lie detector that has always existed---AKS showed how to make it fast.

\subsection{What Was Stopping Progress}
Before 2002, primality testing was split into two incompatible worlds: *fast but possibly wrong* (probabilistic tests like Miller--Rabin) and *always correct but agonizingly slow* (trial division, even elliptic curve methods). Many complexity theorists suspected that deterministic polynomial-time primality testing might be impossible---that PRIMES was not in P at all, or that proving otherwise would require resolving the Riemann Hypothesis. The gap was between an identity everyone knew about and a way to check it efficiently that nobody could find.

\subsection{The Breakthrough}
Agrawal, Kayal, and Saxena's key idea: don't check $(X + a)^n \equiv X^n + a \pmod{n}$ as polynomials of degree $n$. Instead, reduce modulo $X^r - 1$ for a small, carefully chosen $r$, so the polynomials stay short. Then check only a polynomial number of values of $a$. They proved that this---combined with a clever choice of $r$---catches every composite number. The identity itself was old; the novelty was proving a *small* ring suffices.

\subsection{How It Works}

Before 2002, primality testing was split into two incompatible worlds. On one side stood fast algorithms that could be \emph{wrong}. Probabilistic tests like Miller--Rabin flip a coin (figuratively) for each candidate: they run in polynomial time but declare a number ``probably prime'' with a small but nonzero chance of error. For cryptography this was acceptable in practice---you can run the test 100 times and drive the error probability below any threshold---but a theoretical guarantee was impossible. On the other side stood algorithms that were always correct but \emph{slow}. Trial division is exponential: for a $d$-digit number you might need to check up to $10^{d/2}$ divisors. Even the fastest known deterministic methods at the time---elliptic curve primality proving (ECPP)---were only ``practically polynomial'' (heuristically $\tilde{O}((\log n)^4)$) and relied on unproven assumptions about the distribution of elliptic curve group orders.

The key gap was therefore: \textbf{no known deterministic algorithm could prove primality in proven polynomial time, and it was widely believed that this might be impossible.} Many complexity theorists suspected that PRIMES might not be in P at all, or that proving it \emph{would} require resolving deep open questions like the Riemann Hypothesis.

Each previous approach hit a specific wall:

\begin{itemize}
\item \textbf{Fermat's test and its descendants} were undermined by Carmichael numbers---composites like $561 = 3 \cdot 11 \cdot 17$ that satisfy $a^{n-1} \equiv 1 \pmod{n}$ for every $a$ coprime to $n$. No choice of base could reliably distinguish them from primes.
\item \textbf{Miller's deterministic test} (1976) is polynomial-time \emph{assuming} the Extended Riemann Hypothesis (ERH). It reduces the number of bases you need to check from all of $\mathbb{Z}_n^*$ to just $(\log n)^2$, but proving the ERH remains one of the deepest open problems in mathematics. The test is conditional, not unconditional.
\item \textbf{Solovay--Strassen} avoids ERH but is probabilistic: it has a one-in-two error rate per round, and no finite number of rounds eliminates the error entirely.
\item \textbf{Elliptic curve methods} (Goldwasser--Kilian 1986, Atkin--Morain 1993) are deterministic and have expected polynomial running time, but their polynomial-time guarantee is \emph{conditional on heuristics} about the probability that random elliptic curves yield large group orders. No proof that these heuristics always holds was known.
\item \textbf{Special-form tests} (Lucas--Lehmer for Mersenne numbers, Proth's theorem for $k \cdot 2^n + 1$) are fast and unconditional, but they only work for numbers with a very specific algebraic structure. A general-purpose test was missing.
\end{itemize}

The fundamental obstacle was finding a purely algebraic characterization of primality that (a) is \emph{both necessary and sufficient}---not just one direction---and (b) can be checked using computations that scale polynomially with the number of digits. Every prior attempt either got the ``if and only if'' wrong (Fermat) or got it right but could not be checked efficiently without unproven assumptions (Miller, ECPP).

Agrawal, Kayal, and Saxena's breakthrough was to realize that a beautiful identity from abstract algebra could serve as a \emph{perfect} primality test, and---crucially---that it could be made fast.

The starting point is an old fact about prime numbers. If $p$ is prime, then in the ring $\mathbb{Z}/p\mathbb{Z}$:
\[
(X + a)^p \equiv X^p + a \pmod{p}
\]
for every integer $a$. This follows from the binomial theorem and the fact that $\binom{p}{k}$ is divisible by $p$ for $0 < k < p$. In ring-theoretic language, the map $x \mapsto x^p$ is the \emph{Frobenius endomorphism}, and the identity above says that the Frobenius acts as a ring homomorphism---it distributes over addition---if and only if $p$ is prime.

This immediately gives a primality criterion: $n$ is prime if and only if $(X+a)^n \equiv X^n + a \pmod{n}$ as polynomials, for every $a$ coprime to $n$. The ``if'' direction is the hard part, and Agrawal--Kayal--Saxena proved it.

\begin{theorem}[AKS Primality Test, 2002 \cite{agrawal2004}]
There exists a deterministic algorithm that decides whether $n$ is prime in time $\tilde{O}((\log n)^6)$.
\end{theorem}

The novelty was not the identity itself---it was known in algebra for over a century. The novelty was \emph{making it fast}. Na\"ively checking $(X+a)^n \pmod{X^r - 1, n}$ requires multiplying polynomials of degree up to $n$, which is too slow. The key ideas that made it polynomial were:

\begin{enumerate}
\item \textbf{Work modulo $X^r - 1$ for a \emph{small} $r$}. Reducing to degree $< r$ keeps polynomial multiplications manageable (each takes $\tilde{O}(r)$ time). The challenge is choosing $r$ small enough to be fast but large enough to make the test correct.
\item \textbf{Choose $r$ so that the order of $n$ modulo $r$ is large}. Specifically, the algorithm finds $r$ such that $\operatorname{ord}_r(n) > (\log n)^2$. This ensures that if $n$ is composite, the algebraic contradictions propagate far enough to be detected.
\item \textbf{Check only a polynomial number of values of $a$}. Instead of testing every $a$, checking $a = 1, 2, \ldots, \lceil \sqrt{\phi(r)} \log n \rceil$ (where $\phi$ is Euler's totient) suffices. This bounds the total work.
\end{enumerate}

Here is the intuition in plain language. To test whether $n$ is prime:

\begin{enumerate}
\item Pick a small integer $r$ (chosen so that $n$ has large multiplicative order modulo $r$).
\item In the polynomial ring $\mathbb{Z}_n[X] / (X^r - 1)$, compute $(X + a)^n$ for several values of $a$.
\item Check whether the result equals $X^n + a$ in that ring.
\item If it fails for any $a$, then $n$ is composite. If it passes for all tested values of $a$, then $n$ is prime.
\end{enumerate}

Think of it this way: \emph{composite numbers cannot consistently fake the Frobenius endomorphism}. When you expand $(X+a)^n$ and reduce modulo both $n$ and $X^r - 1$, a composite $n$ will ``leak''---the coefficients will not match $X^n + a$. The genius of AKS was proving that checking a polynomial number of $a$-values in a ring of polynomial degree suffices to catch every composite.

\subsection{Complexity and Improvements}

The original AKS paper ran in $\tilde{O}((\log n)^{12})$. Subsequent improvements dramatically reduced the running time:

\begin{itemize}
\item Bernstein (2003): practical optimizations, $\tilde{O}((\log n)^{12})$ with smaller constants.
\item Lenstra and Pomerance (2011): improved the theoretical bound to $\tilde{O}((\log n)^3)$, which is close to optimal \cite{lenstra2011}.
\item Cohen and Lenstra (2003): various structural optimizations \cite{bernstein2006}.
\end{itemize}

The algorithm's correctness proof uses deep algebraic number theory:

\begin{enumerate}
\item \textbf{Cyclotomic fields}: The proof studies splitting of primes in the $r$-th cyclotomic field $\mathbb{Q}(\zeta_r)$, using the fact that the Frobenius element controls how primes factor in extensions.
\item \textbf{Order arguments}: The choice of $r$ ensures that if $n$ is composite and passes the polynomial tests, then $n$ must have a very specific algebraic structure---it must ``look like'' a prime in many cyclotomic extensions simultaneously.
\item \textbf{Contradiction via irreducibility}: The proof shows that this structure is impossible for composite numbers: if $n$ is composite, the polynomial $X^r - 1$ would need to factor in a way that contradicts known results about cyclotomic polynomials.
\end{enumerate}

\begin{highlightbox}
The AKS algorithm was the first \emph{general} deterministic polynomial-time primality test. It settled a decades-old open problem and proved unconditionally that PRIMES $\in$ P. While not the fastest in practice (probabilistic tests are still preferred for real-world applications), it is theoretically fundamental: it showed that primality is a property that can be checked in polynomial time, without any unproven assumptions. The proof introduced a novel bridge between computational complexity and algebraic number theory that continues to inspire research.
\end{highlightbox}

\subsection{After the Proof}

The AKS result is completely settled: PRIMES $\in$ P is a theorem, and the question of whether deterministic polynomial-time primality testing is possible has a definitive answer. The proof is unconditional---it does not rely on the Riemann Hypothesis or any other unproven assumption. In that sense, the theoretical question is closed.

However, the AKS algorithm is not practical. It is far too slow compared to probabilistic tests like Miller--Rabin, which are both faster and simple enough to implement. In real-world cryptography, nobody uses AKS. Its importance is theoretical: it settled a fundamental question about the boundary between efficient and inefficient computation. The proof also introduced new connections between algebraic number theory (cyclotomic fields, order of elements in multiplicative groups) and computational complexity that continue to inspire research.

On the quantitative side, room for improvement remains. The original AKS algorithm ran in $\tilde{O}((\log n)^{12})$ time, which was quickly reduced to $\tilde{O}((\log n)^6)$ and eventually to $\tilde{O}((\log n)^3)$ by Lenstra and Pomerance in 2011. Whether the exponent can be reduced further---perhaps to $\tilde{O}((\log n)^2)$ or even $\tilde{O}(\log n)$---is an open question. The Lenstra--Pomerance bound is conjectured to be close to optimal, but a matching lower bound is not known. Beyond speed, there is also the question of whether the underlying ideas can be used to construct other efficient deterministic algorithms for related algebraic problems.

\section{Impact}

\begin{itemize}
\item \textbf{Complexity theory}: Proved that PRIMES $\in$ P, resolving a major open problem.
\item \textbf{Cryptography}: While not used in practice due to slower performance, it provides a theoretical foundation for primality testing.
\item \textbf{Mathematical techniques}: The proof introduced novel use of cyclotomic polynomials and order arguments that have influenced other areas.
\end{itemize}

\section{Computational Verification}

The AKS primality test can be implemented and verified computationally. The following Python code demonstrates the core ideas.

\subsection{The AKS Criterion in Python}

The simplest version of the AKS test checks whether $(X+a)^n \equiv X^n + a \pmod{n}$ as polynomials.

\begin{verbatim}
def aks_criterion(n, a=2):
    """Check if n passes the AKS polynomial test for base a."""
    # Expand (X + a)^n mod n as a polynomial
    # Use the binomial theorem: coefficient of X^k is C(n,k) * a^(n-k)
    coeffs = []
    binom = 1  # C(n, 0) = 1
    for k in range(n + 1):
        coeff = (binom * pow(a, n - k, n)) % n
        coeffs.append(coeff)
        # Update binomial coefficient: C(n, k+1) = C(n, k) * (n - k) / (k + 1)
        binom = (binom * (n - k) // (k + 1)) if k < n else 0

    # Check: should equal X^n + a, i.e., coeff[0] = a mod n,
    # coeff[n] = 1 mod n, all others = 0 mod n
    for k in range(1, n):
        if coeffs[k] % n != 0:
            return False  # composite
    return coeffs[0] % n == a % n and coeffs[n] % n == 1

# Test on small numbers
for n in range(2, 50):
    if aks_criterion(n):
        assert all(n % p != 0 for p in range(2, int(n**0.5)+1) or n == p)
        print(f"{n}: prime (AKS confirms)")
    else:
        print(f"{n}: composite")
\end{verbatim}

\subsection{Carmichael Numbers: Fermat Test Fails}

Carmichael numbers pass the Fermat test for every base coprime to $n$, but fail the AKS test.

\begin{verbatim}
def fermat_test(n, a):
    """Fermat primality test: does a^(n-1) = 1 mod n?"""
    return pow(a, n - 1, n) == 1

carmichael = [561, 1105, 1729, 2465, 2821, 6601]
for n in carmichael:
    # Fermat test passes for all coprime bases
    fermat_results = [fermat_test(n, a) for a in range(2, n) if gcd(a, n) == 1]
    all_pass = all(fermat_results)
    # AKS test correctly identifies composite
    aks_result = aks_criterion(n)
    print(f"{n}: Fermat all-pass={all_pass}, AKS prime={aks_result}")
\end{verbatim}

\subsection{Timing Comparison}

Compare AKS with trial division and Miller--Rabin on small inputs.

\begin{verbatim}
import time

def trial_division(n):
    if n < 2: return False
    for i in range(2, int(n**0.5) + 1):
        if n % i == 0: return False
    return True

# Time on a moderately large prime
n = 104729  # 10000th prime
start = time.time()
aks_criterion(n)
aks_time = time.time() - start

start = time.time()
trial_division(n)
td_time = time.time() - start

print(f"n = {n}")
print(f"AKS time:              {aks_time:.4f}s")
print(f"Trial division time:   {td_time:.4f}s")
\end{verbatim}

\begin{highlightbox}
Students should experiment with these scripts to observe: (1)~how Carmichael numbers fool the Fermat test but not AKS, (2)~how AKS scales compared to trial division, and (3)~the effect of choosing different values of $r$ in the full AKS algorithm.
\end{highlightbox}

\section{Exercises}

\begin{exercise}[Basic]
Verify that $(X + a)^5 \equiv X^5 + a \pmod{5}$ for any integer $a$. Check that the binomial coefficients $\binom{5}{1}, \binom{5}{2}, \binom{5}{3}, \binom{5}{4}$ are all divisible by 5.
\end{exercise}

\begin{exercise}[Intermediate]
Show that 561 = 3 $\cdot$ 11 $\cdot$ 17 is a Carmichael number: $a^{560} \equiv 1 \pmod{561}$ for all $a$ coprime to 561.
\end{exercise}

\begin{exercise}[Challenge]
Explain why the polynomial congruence $(X+a)^n \equiv X^n + a \pmod{n}$ is a valid primality test. Why doesn't this work directly as a fast algorithm?
\end{exercise}

\begin{exercise}
The AKS algorithm runs in $\tilde{O}((\log n)^6)$. If $n$ has $d$ digits, how does this compare to trial division? What is the practical implication?
\end{exercise}


\begin{exercise}[Computational]
Write a Python/SageMath script to explore a concept from this chapter. See the appendix for starter code.
\end{exercise}

\section{Timeline}

\begin{tabularx}{\linewidth}{lX}
\toprule
\textbf{Year} & \textbf{Event} \\ \midrule
1976 & Miller proposes a deterministic primality test (conditional on GRH) \\ 
1976 & Rabin develops the Miller--Rabin probabilistic test \\ 
1980s & Various probabilistic tests (Solovay--Strassen, Fermat test) \\ 
2002 & Agrawal, Kayal, Saxena publish the AKS algorithm \cite{agrawal2004} \\ 
2003 & Cohen and Lenstra optimize the algorithm \cite{bernstein2006} \\ 
2011 & Lenstra and Pomerance achieve $\tilde{O}((\log n)^3)$ \cite{lenstra2011} \\ \bottomrule
\end{tabularx}

\section{Citations}

\begin{enumerate}
\item Agrawal, M., Kayal, N., \& Saxena, N. (2004). ``PRIMES is in P.'' Annals of Mathematics, 160(2), 781--793.
\item Lenstra, H. W., \& Pomerance, C. (2011). ``Highly robust error-correcting polynomial primality testing.'' Journal of the ACM, 58(1), Article 6.
\item Bernstein, D. J. (2006). ``Speeding up the AKS primality test.'' In Algorithmic Number Theory.
\end{enumerate}
